\section{Evaluation}
We evaluate Paper Circle along three axes: (i) retrieval effectiveness under different configurations, (ii) stability and reproducibility of rankings across steps, and (iii) the utility of diversity-aware postprocessing for surfacing non-redundant results. Paper Circle provides built-in evaluation metrics but does not enforce a fixed benchmark dataset. When a ground-truth paper title or identifier is provided, the system computes Mean Reciprocal Rank (MRR), Recall@K, Precision@K, and hit rates. These metrics are computed per step and stored in \texttt{retrieval\_metrics.json} for longitudinal tracking.

As a minimal illustrative scenario, consider a known target paper in the local corpus: the pipeline is run once using offline retrieval and once using online sources. The resulting MRR and Recall@K values allow direct comparison of configuration impact, while repeated runs confirm stable rankings when deterministic scoring is enabled. Although lightweight, this framing aligns evaluation with discovery goals rather than task-specific QA benchmarks.

For batch evaluation, a parallel benchmarking utility executes multiple queries concurrently and aggregates mean metrics and timing statistics. This supports lightweight comparisons between search configurations (offline vs. online, BM25 vs. semantic, with or without reranking) without requiring external tooling.
